% !PRINTED 414
\weberpage{414}
%<<CONT
werden können, dass die in (10) vorkommende Constante $\varepsilon$ \gesperrt{immer}
eine dritte Einheitswurzel ist. Wenn $\pi$ zu den $\alpha$ gehört, so ist
$\varepsilon$ zugleich eine zweite Einheitswurzel und muss also $=1$ sein.
Daraus folgern wir den Satz:

2. \gesperrt{Die Function $\vartheta$ bleibt bei den Permutationen der
Vierergruppe $A$ ungeändert.}

Die dritte Potenz von $\vartheta$ gestattet die Permutationen der
alternirenden Gruppe und $\vartheta$ ist daher die Wurzel einer reinen
cubischen Gleichung $\vartheta^3 = z$. Dies ist also die Resolvente
$\Phi(\vartheta,z) = 0$, die in diesem Falle eine \gesperrt{Partialresolvente} ist
(Bd. I, \S.\,161).

Aus diesen Betrachtungen ergiebt sich nun ohne
Schwierigkeit der folgende Satz, dessen Beweis das Ziel dieser ganzen
Betrachtungen ist:

3. \gesperrt{Eine allgemeine Gleichung von höherem als
4\textsuperscript{ten} Grade hat keine Resolventen mit einem Parameter.}

Denn ist $n$ grösser als 4, so können wir, wenn wir $\vartheta$ als
Function von je vieren der Variablen $x$ betrachten, den Satz 2. anwenden.
Es folgt dann, dass $\vartheta$ ungeändert bleibt, wenn unter den
Variablen irgend ein Paar von Transpositionen vorgenommen wird. Da man
nun aus Transpositionspaaren die ganze alternirende Gruppe zusammensetzen
kann (Bd. I, \S.\,153, 8.), so folgt, dass $\vartheta$ überhaupt durch
die alternirende Gruppe ungeändert bleibt, also eine alternirende oder
eine symmetrische Function ist. $\Phi(\vartheta,z) = 0$ reducirt sich auf
die Identität $\vartheta = z$ und ist keine Resolvente.

Es sei noch bemerkt, dass, wenn man an Stelle der alternirenden Gruppe
die symmetrische treten lässt, dieser Satz a fortiori gilt.

Den Satz, dessen Beweis wir hier entwickelt haben, hat zuerst
\gesperrt{Kronecker} ausgesprochen, ohne einen Beweis dafür zu
veröffentlichen. Der erste Beweis ist von \gesperrt{F. Klein} gegeben, der
dann von \gesperrt{Gordan} noch vereinfacht
ist\footnote{\gesperrt{F. Klein}, Vorlesungen über das Ikosaëder. Der
Beweis von \gesperrt{Gordan}, dem unsere Darstellung in den Grundzügen
folgt, findet sich in Bd. 29 der mathematischen Annalen (1887).}.
